\section{Основная часть}

\subsection{Лабораторная работа № 1}

\subsubsection{Задача}

В ходе лабораторной работы необходимо:

\begin{enumerate}
    \item создать представление, содержащее результат группировки данных с
    двумя агрегатными функциями \mintinline{sql}{COUNT};
    \item составить простой запрос к созданному представлению с условием
    \mintinline{sql}{WHERE};
    \item составить запрос, в котором представление используется совместно с
    операциями соединения \mintinline{sql}{JOIN}.
\end{enumerate}

Представление должно хранить число форматов и градаций насыщенности каждого
начертания. Простой запрос должен определять число начертаний, поддерживающих
больше $A$ форматов. Запрос с соединениями должен находить пять категорий с
наибольшим охватом форматов.

В формулах используются обозначения: $T$~--- таблица
\mintinline{text}{Typeface}, $FF$~--- таблица
\mintinline{text}{Font_family}, $C$~--- таблица
\mintinline{text}{Category}, $V$~--- созданное представление. Символ
$\bowtie$ обозначает соединение, $\sigma$~--- выборку, $\pi$~--- проекцию,
$\gamma$~--- группировку с агрегированием.

\subsubsection{Создание представления}

Одно начертание может храниться в таблице \mintinline{text}{Typeface} в
нескольких форматах и градациях насыщенности. Чтобы не выполнять одинаковый
подсчёт в каждом запросе, создаётся отдельное представление.

В нём для каждого начертания хранятся название, идентификатор гарнитуры, число
форматов и число градаций насыщенности. Определение представления приведено в
листинге~\ref{lst:create-view}.

\codeListing[sql]{../practice-01/src/create.sql}{Создание представления}{\label{lst:create-view}}

Агрегатная функция \mintinline{sql}{COUNT(DISTINCT ...)} подсчитывает
количество различных значений. Поэтому повторяющиеся форматы и градации
насыщенности учитываются только один раз.

Создание представления описывается формулой

\[
V = \gamma_{\substack{T.id\_font\_family,T.name;\\
\operatorname{count\_distinct}(T.id\_format) \rightarrow format\_count,\\
\operatorname{count\_distinct}(T.id\_weight) \rightarrow weight\_count}}(T).
\]

\subsubsection{Проверка представления}

Для проверки из представления выбираются первые десять начертаний. Запрос
приведён в листинге~\ref{lst:query-01}.

\codeListing[sql]{../practice-01/src/query-01.sql}{Проверка содержимого представления}{\label{lst:query-01}}

В реляционная формула имеет вид

\[
R_1 = \pi_{\substack{V.id\_font\_family,V.typeface\_name,\\
V.format\_count,V.weight\_count}}(V).
\]

План и результат выполнения показаны на
рисунках~\ref{fig:query-01-plan} и~\ref{fig:query-01-result}.

\img[0.55]{img/query-01-plan.png}{План выполнения запроса 1}{\label{fig:query-01-plan}}

\img[0.8]{img/query-01-result.png}{Результат выполнения запроса 1}{\label{fig:query-01-result}}

В результате для каждого начертания получена одна строка с двумя рассчитанными
значениями: \mintinline{sql}{format_count} и
\mintinline{sql}{weight_count}.

\clearpage
\subsubsection{Запрос с условием}

В запросе значение $A$ принято равным 2. Сначала условие
\mintinline{sql}{format_count > 2} оставляет подходящие начертания, затем
\mintinline{sql}{COUNT} подсчитывает их. Запрос приведён в
листинге~\ref{lst:query-02}.

\codeListing[sql]{../practice-01/src/query-02.sql}{Подсчёт начертаний с числом форматов больше
$A$}{\label{lst:query-02}}

Запрос описывается формулой

\[
R_2 = \gamma_{\operatorname{count}(V.id\_font\_family)
\rightarrow typeface\_count}
\left(\sigma_{V.format\_count>A}(V)\right).
\]

План и результат выполнения показаны на
рисунках~\ref{fig:query-02-plan} и~\ref{fig:query-02-result}.

\img[0.42]{img/query-02-plan.png}{План выполнения запроса 2}{\label{fig:query-02-plan}}

\img[0.7]{img/query-02-result.png}{Результат выполнения запроса 2}{\label{fig:query-02-result}}

Условию соответствует 39 начертаний.

\subsubsection{Охват форматов по категориям}

Представление содержит идентификатор гарнитуры, поэтому его можно соединить с
таблицами \mintinline{text}{Font_family} и \mintinline{text}{Category}. После
соединения значения \mintinline{sql}{format_count} суммируются отдельно для
каждой категории. Запрос приведён в листинге~\ref{lst:query-03}.

\codeListing[sql]{../practice-01/src/query-03.sql}{Первые пять категорий по охвату
поддерживаемых форматов}{\label{lst:query-03}}

Агрегирование по категориям и выбор первых пяти строк описываются формулами

\[
\begin{aligned}
Q &= \gamma_{\substack{C.id\_category,C.name;\\
\operatorname{sum}(V.format\_count) \rightarrow format\_count}}
(V \bowtie FF \bowtie C),\\
R_3 &= \operatorname{top}_{5}^{format\_count\downarrow,\,C.name}(Q).
\end{aligned}
\]

Представление уже содержит число форматов каждого начертания, поэтому повторно
считать их в запросе не требуется. План и результат выполнения показаны на
рисунках~\ref{fig:query-03-plan} и~\ref{fig:query-03-result}.

\img[0.75]{img/query-03-plan.png}{План выполнения запроса 3}{\label{fig:query-03-plan}}

\img[0.75]{img/query-03-result.png}{Результат выполнения запроса 3}{\label{fig:query-03-result}}

После группировки категории сортируются по убыванию охвата. Оператор
\mintinline{sql}{LIMIT 5} оставляет первые пять строк. Наибольший охват в
полученном результате имеет категория \mintinline{text}{Neo-grotesque}.
